\lesson{4}{Oct 18 2021 Mon (10:55:28)}{Graphing Quadratics}{Unit 2}

\subsubsection*{Parts of a Parabola}

Parabolas consist of a:

\begin{itemize}
    \item Vertex
    \item Axis of Symmetry
    \item Domain
    \item Range
    \item X-Y Intercepts (Sometimes though)
\end{itemize}

\paragraph{Vertex} Here is how you find the vertex of a parabola:

\begin{itemize}
    \item \textbf{Step One: $x$-Coordinate} Set the expression inside the parentheses equal to $0$ and solve for the value of $x$
    \item \textbf{Step Two: $y$-Coordinate} Substitute the value of $x$ from step one into the original equation and solve for $y$
\end{itemize}

Let's try an example:

\begin{example}[Given the equation $f(x) = -2(x - 5)^2 + 8$, find the $(x,y)$ coordinate of the vertex]
    \begin{align}
        f(x) &= -2(x - 5)^2 + 8 \\
             &= -2([x - 5 = 0])^2 + 8
             &= 5 - 5 = 0 \\
             &= x = 5
    \end{align}
    
    Now, we're going to solve for the $y$ coordinate
    
    \begin{align}
        f(x) &= -2(5)(5 - 5)^2 + 8 \\
             &= -2(5 - 5)^2 + 8 \\
             &= -2(0)^2 + 8 \\
             &= -2(0) + 8 \\
             &= y = 8 \\
    \end{align}
    
    Now, we know that the vertex of the parabola is $(5, 8)$.
\end{example}

\newpage

\subsubsection*{Elements of the Parabola}

\begin{marginfigure}
    \centering
    \incfig{parabola}
    \sidecaption{Here is an example of what a parabola looks like.}
    \label{fig:parabola}
\end{marginfigure}

\paragraph{Line of Symmetry} The \textbf{axis of symmetry} is in the line
that divides the parabola into two equal parts where each part is a
mirror reflection of the other.

As you can see, the axis of symmetry is along the line $x = -3$.

To find the equation of the axis of symmetry algebraically, set the expression inside the parentheses equal to $0$ and solve for $x$.

\paragraph{Domain and Range} Any number can be substituted for the variable $x$
to produce a unique $y$ value. Therefore, the domain is \textbf{"All Real Numbers}.

This parabola opens up since the value is positive ($+$). Because the parabola opens up, the vertex represents the minimum point on the graph.
All of the $y$ values on the parabola must be above the minimum $y-coordinate$ of $-1$, which makes the range of the parabola this:

\begin{align}
    y \geq -1
\end{align}

\paragraph{Intercepts} Substituting $0$ for $x$ and solving for $y$ will result in the $y-intercept$:

\begin{align}
    f(g) &= (x + 3)^2 - 1 \\
         &= (0 + 3)^2 - 1 \\
         &= (3)^2 - 1 \\
         &= 9 - 1 \\
         &= 8
\end{align}

There you go, you have your $y-intercept$. To determine the $x-intercept$, you just plug in the $y-intercept$:

\begin{align}
    f(9) &= (9 + 3)^2 - 1 \\
         &= (3)^2 - 1 \\
         &= 9 - 1 \\
         &= 8
\end{align}

\subsubsection*{Different Quadratic Equation forms}

While identifying the:

\begin{enumerate}
    \item Vertex
    \item Axis of Symmetry
    \item Domain
    \item Range
    \item Intercepts
\end{enumerate}

is most easily done using the vertex form of a quadratic equation, most of the time, quadratic equations are written in standard form.

Here is the standard form of a quadratic equation:

\begin{align}
    f(x) = ax^2 + bx + c
\end{align}

Now, we can move that into vertex form, here is an example:

\begin{align}
    f(x) &= -2x(x - 5)^2 + 8 \\
         &= -2x(x - 5)(x - 5) + 8 \\
         &= -2x(x^2 - 5x - 5x + 25) + 8 \\
         &= -2x(x^2 - 10x + 25) + 8 \\
         &= -2x^2 + 20x - 50 + 8 \\
         &= -2x^2 + 20x - 42 \\
\end{align}

But, it's not that easy to identify the vertex and axis, so let's us \textbf{Axis of Symmetry} equation, which is:

\begin{align}
    f(x) = -\frac{b}{2a}
\end{align}

Now, let's try to convert an equation in the standard form to this form:

\begin{align}
    f(x) &= -2x^2 + 20x - 42 \\
         &= -\frac{20}{2(-2)} \\
         &= -\frac{20}{-4} \\
         &= x = 5
\end{align}

which gives you the axis of symmetry for the $x-axis$. Now, just substitute $5$ for $x$ to find the $y-coordinate$:

\begin{align}
    f(x) &= -2x^2 + 20x - 42 \\
         &= -2(5)^2 + 20(5) - 42 \\
         &= -2(25) + 100 - 42 \\
         &= -50 + 100 - 42 \\
         &= 8 \\
         &= (5, 8)
\end{align}

\subsubsection*{Working our way up}

\begin{marginfigure}
    \centering
    \incfig{parabola}
    \sidecaption{Here is an example of what a parabola looks like.}
    \label{fig:parabola}
\end{marginfigure}

Now, I have discussed how you can find key features within the equation, but let's try and make our equation from the graph.
Taking the image of the parabola, let's try and see if we can create the equation from the graph:

Let's take the vertex $(-3, -1)$ and a random point, $(0, 8)$. Here is how it works:

\begin{align}
    f(x) &= a(x - h)^2 + k \\
         &= a(x + 3)^2 - 1 \\
         &= 8 = a(0 + 3)^2 - 1 \\
    8 &= a(0 + 3)^2 - 1 \\
       &= a(3)^2 - 1 \\
       &= 9a - 1 \\
       &= 9 = 9a \\
       &= \frac{9}{9} = \frac{9}{9a} \\
       &= 1
\end{align}

Now, let's put it all together:

\begin{align}
    a &= 1, h = -3, k = -1 \\
    f(x) &= a(x - h)^2 + k \\
         &= (x + 3)^2 - 1 \\
         &= (x + 3)(x + 3) - 1 \\
         &= (x^2 + 3x + 3x + 9) - 1 \\
         &= (x^2 + 6x + 9) - 1 \\
         &= x^2 + 6x + 8 \\
\end{align}

\newpage
